\section{Problem 4: Partial Resolution and a Program Toward the General Case}\label{sec:q4}

Let $n\ge 2$. For monic degree-$n$ polynomials
\[
p(x)=\sum_{k=0}^n a_k x^{n-k},\qquad q(x)=\sum_{k=0}^n b_k x^{n-k},\qquad a_0=b_0=1,
\]
define
\[
(p\boxplus_n q)(x)=\sum_{k=0}^n c_k x^{n-k},
\qquad
c_k=\sum_{i+j=k}\frac{(n-i)!(n-j)!}{n!(n-k)!}\,a_i b_j.
\]
For real-rooted $p$ with roots $\lambda_1,\dots,\lambda_n$, define
\[
\Phi_n(p):=\sum_{i=1}^n\left(\sum_{j\ne i}\frac{1}{\lambda_i-\lambda_j}\right)^2,
\]
and set $\Phi_n(p)=\infty$ if $p$ has a multiple root.

The target inequality is
\[
\frac{1}{\Phi_n(p\boxplus_n q)}\ge \frac{1}{\Phi_n(p)}+\frac{1}{\Phi_n(q)}.
\tag{$\star$}\label{eq:q4-star}
\]

\begin{remark}[Strict Status]
This section gives complete proofs of several structural identities and special cases, and records a concrete route for further progress. The all-$n$ statement \eqref{eq:q4-star} is not proved here.
\end{remark}

\begin{lemma}[Pairwise inverse-square identity]\label{lem:q4-pairwise}
If $p$ has simple roots $\lambda_1,\dots,\lambda_n$, then
\[
\Phi_n(p)=\sum_{i\ne j}\frac{1}{(\lambda_i-\lambda_j)^2}
=2\sum_{1\le i<j\le n}\frac{1}{(\lambda_i-\lambda_j)^2}.
\]
\end{lemma}

\begin{proof}
Set $S_i:=\sum_{j\ne i}(\lambda_i-\lambda_j)^{-1}$. Expanding $\sum_i S_i^2$ gives the pairwise inverse-square term plus a triple sum.
For distinct $a,b,c$, one has
\[
\frac{1}{(a-b)(a-c)}+\frac{1}{(b-a)(b-c)}+\frac{1}{(c-a)(c-b)}=0,
\]
so triple contributions cancel by grouping unordered triples of indices.
\end{proof}

\begin{proposition}[Exact equality for $n=2$]\label{prop:q4-n2}
For all monic real-rooted quadratics $p,q$,
\[
\frac{1}{\Phi_2(p\boxplus_2 q)}=\frac{1}{\Phi_2(p)}+\frac{1}{\Phi_2(q)}.
\]
\end{proposition}

\begin{proof}
Write
\[
p=x^2+a_1x+a_2,\qquad q=x^2+b_1x+b_2,\qquad r=p\boxplus_2 q=x^2+c_1x+c_2.
\]
From the definition,
\[
c_1=a_1+b_1,\qquad c_2=a_2+b_2+\tfrac12 a_1b_1.
\]
Hence
\[
\Delta_r:=c_1^2-4c_2=(a_1^2-4a_2)+(b_1^2-4b_2)=\Delta_p+\Delta_q.
\]
For a monic quadratic with distinct real roots and gap $d$, one has $\Delta=d^2$ and by Lemma~\ref{lem:q4-pairwise},
\[
\Phi_2=\frac{2}{d^2}=\frac{2}{\Delta},\qquad \frac{1}{\Phi_2}=\frac{\Delta}{2}.
\]
Thus
\[
\frac{1}{\Phi_2(r)}=\frac{\Delta_r}{2}=\frac{\Delta_p}{2}+\frac{\Delta_q}{2}.
\]
When a polynomial has a repeated root, $\Phi_2=\infty$ and the same identity holds with the convention $1/\infty=0$.
\end{proof}

\begin{proposition}[Full inequality for $n=3$]\label{prop:q4-n3}
For all monic real-rooted cubics $p,q$,
\[
\frac{1}{\Phi_3(p\boxplus_3 q)}\ge \frac{1}{\Phi_3(p)}+\frac{1}{\Phi_3(q)}.
\]
\end{proposition}

\begin{proof}
Step 1 (normalization by translation).
For $a,b\in\mathbb R$, let
\[
\tau_a p(x):=p(x-a),\qquad \tau_b q(x):=q(x-b).
\]
Using the permutation-average model of $\boxplus_3$ on roots, one has
\[
(\tau_a p)\boxplus_3(\tau_b q)=\tau_{a+b}(p\boxplus_3 q).
\]
Since $\Phi_3$ depends only on root differences, it is translation-invariant:
\[
\Phi_3(\tau_c r)=\Phi_3(r).
\]
Hence it suffices to prove the inequality after translating $p,q$ so that each has root mean $0$.

Step 2 (centered form).
Write
\[
p(x)=x^3+A_1x+B_1,\qquad q(x)=x^3+A_2x+B_2
\]
with $A_1,A_2<0$ for nontrivial real-rooted cubics.
From the coefficient definition of $\boxplus_3$:
\[
p\boxplus_3 q=x^3+(A_1+A_2)x+(B_1+B_2).
\]

Step 3 (closed formula for $1/\Phi_3$).
For a centered monic cubic $r(x)=x^3+Ax+B$ with simple real roots:
\[
\Delta(r)=-4A^3-27B^2,\qquad
\sum_{i<j}(\lambda_i-\lambda_j)^2=-6A.
\]
For three numbers, if $a,b,c$ are the three squared pairwise differences, then
$ab+ac+bc=(a+b+c)^2/4$. Therefore
\[
\Phi_3(r)=2\sum_{i<j}\frac{1}{(\lambda_i-\lambda_j)^2}
      =\frac{\bigl(\sum_{i<j}(\lambda_i-\lambda_j)^2\bigr)^2}{2\,\Delta(r)}
      =\frac{18A^2}{\Delta(r)}.
\]
So
\[
\frac{1}{\Phi_3(r)}=\frac{\Delta(r)}{18A^2}
=-\frac{2A}{9}-\frac{3}{2}\frac{B^2}{A^2}.
\]

Apply this to $p,q,r:=p\boxplus_3 q$:
\[
\frac{1}{\Phi_3(r)}-\frac{1}{\Phi_3(p)}-\frac{1}{\Phi_3(q)}
=-\frac{3}{2}\left(
\frac{(B_1+B_2)^2}{(A_1+A_2)^2}
-\frac{B_1^2}{A_1^2}
-\frac{B_2^2}{A_2^2}\right),
\]
because the linear $A$-terms cancel.

Set $u:=-A_1>0$, $v:=-A_2>0$, $x:=B_1/A_1$, $y:=B_2/A_2$. Then
\[
\frac{(B_1+B_2)^2}{(A_1+A_2)^2}=\left(\frac{ux+vy}{u+v}\right)^2.
\]
By Cauchy (or Jensen for $t\mapsto t^2$),
\[
\left(\frac{ux+vy}{u+v}\right)^2\le \frac{ux^2+vy^2}{u+v}\le x^2+y^2
=\frac{B_1^2}{A_1^2}+\frac{B_2^2}{A_2^2}.
\]
Hence the bracket is $\le 0$, so the displayed difference is $\ge 0$.

If one polynomial has multiple roots, the convention $1/\infty=0$ and continuity from the simple-root regime give the same inequality. Therefore the claim holds for all real-rooted cubics.
\end{proof}

\begin{proposition}[Degenerate shift case]\label{prop:q4-shift}
If $p(x)=(x-a)^n$, then
\[
p\boxplus_n q(x)=q(x-a),
\]
and
\[
\frac{1}{\Phi_n(p\boxplus_n q)}=\frac{1}{\Phi_n(p)}+\frac{1}{\Phi_n(q)}.
\]
\end{proposition}

\begin{proof}
Using the standard permutation-average model of symmetric additive convolution,
\[
(p\boxplus_n q)(x)=\frac1{n!}\sum_{\pi\in S_n}\prod_{r=1}^n(x-\lambda_r-\mu_{\pi(r)}),
\]
if all $\lambda_r=a$, each summand equals $\prod_r(x-a-\mu_r)=q(x-a)$. Root gaps are translation-invariant, so
$\Phi_n(q(\cdot-a))=\Phi_n(q)$. Since $\Phi_n((x-a)^n)=\infty$, the equality follows.
\end{proof}

\begin{proposition}[Variance additivity under $\boxplus_n$]\label{prop:q4-var}
Let $m(p)$ be the empirical mean of the roots of $p$ and $\operatorname{Var}(p)$ the empirical variance. Then
\[
\operatorname{Var}(p\boxplus_n q)=\operatorname{Var}(p)+\operatorname{Var}(q).
\]
\end{proposition}

\begin{proof}
For $p=x^n+a_1x^{n-1}+a_2x^{n-2}+\cdots$,
\[
\sum_i\lambda_i=-a_1,\qquad \sum_i\lambda_i^2=a_1^2-2a_2.
\]
For $r=p\boxplus_n q$, one has
\[
c_1=a_1+b_1,\qquad c_2=a_2+b_2+\frac{n-1}{n}a_1b_1.
\]
Hence
\[
\sum_i\gamma_i^2=c_1^2-2c_2=(a_1^2-2a_2)+(b_1^2-2b_2)+\frac{2}{n}a_1b_1,
\]
while
\[
\frac1n\Big(\sum_i\gamma_i\Big)^2=\frac1n(a_1+b_1)^2.
\]
The mixed term cancels exactly in
\[
\operatorname{Var}(r)=\frac1n\sum_i\gamma_i^2-\frac1{n^2}\Big(\sum_i\gamma_i\Big)^2,
\]
yielding additivity.
\end{proof}

\begin{lemma}[de Bruijn-type identity for polynomial heat flow]\label{lem:q4-debruijn}
Let $p_t=e^{t\partial_{xx}}p$ and assume $p_t$ has simple real roots $\lambda_1(t),\dots,\lambda_n(t)$ on an interval of $t$. Then
\[
\frac{d}{dt}\log\operatorname{Disc}(p_t)=-4\,\Phi_n(p_t),
\]
where $\operatorname{Disc}(p_t)=\prod_{i<j}(\lambda_i(t)-\lambda_j(t))^2$.
\end{lemma}

\begin{proof}
Differentiate $p_t(\lambda_i(t))=0$:
\[
0=\partial_t p_t(\lambda_i)+p_t'(\lambda_i)\lambda_i'=p_t''(\lambda_i)+p_t'(\lambda_i)\lambda_i',
\]
so
\[
\lambda_i'=-\frac{p_t''(\lambda_i)}{p_t'(\lambda_i)}=-2\sum_{j\ne i}\frac{1}{\lambda_i-\lambda_j}.
\]
Also,
\[
\log\operatorname{Disc}(p_t)=2\sum_{i<j}\log|\lambda_i-\lambda_j|.
\]
Differentiate and substitute the formula for $\lambda_i'$ to get
\[
\frac{d}{dt}\log\operatorname{Disc}(p_t)
=\sum_i \lambda_i'\Big(2\sum_{j\ne i}\frac{1}{\lambda_i-\lambda_j}\Big)
=-4\sum_i\Big(\sum_{j\ne i}\frac{1}{\lambda_i-\lambda_j}\Big)^2
=-4\Phi_n(p_t).
\]
\end{proof}

\begin{lemma}[Differential-operator representation of $\boxplus_n$]\label{lem:q4-diffrepr}
For any monic degree-$n$ polynomials $p,q$,
\[
(p\boxplus_n q)(x)=\frac{1}{n!}\sum_{k=0}^{n} p^{(k)}(x)\,q^{(n-k)}(0).
\]
\end{lemma}

\begin{proof}
Write
\[
p(x)=\sum_{i=0}^n a_i x^{n-i},\qquad q(x)=\sum_{j=0}^n b_j x^{n-j}.
\]
For fixed $k$,
\[
p^{(k)}(x)=\sum_{i=0}^{n-k} a_i\frac{(n-i)!}{(n-i-k)!}x^{n-i-k},
\qquad
q^{(n-k)}(0)=b_k (n-k)!.
\]
Hence
\[
\frac1{n!}\sum_{k=0}^{n} p^{(k)}(x)q^{(n-k)}(0)
=
\sum_{k=0}^{n}\sum_{i=0}^{n-k}
a_i b_k\frac{(n-i)!(n-k)!}{n!(n-i-k)!}x^{n-i-k}.
\]
Set $m=i+k$. The coefficient of $x^{n-m}$ is
\[
\sum_{i+k=m}\frac{(n-i)!(n-k)!}{n!(n-m)!}a_i b_k,
\]
which is exactly the defining coefficient formula of $p\boxplus_n q$.
\end{proof}

\begin{lemma}[Transform factorization and elementary-step decomposition]\label{lem:q4-transform-factor}
Define
\[
\mathcal T_n\!\left(\sum_{k=0}^{n}a_k x^{n-k}\right):=
\sum_{k=0}^{n}(n-k)!\,a_k z^k\in \mathbb R[z]/(z^{n+1}).
\]
Then:
\begin{enumerate}
\item
\[
\mathcal T_n(p\boxplus_n q)=\frac1{n!}\,\mathcal T_n(p)\mathcal T_n(q).
\]
\item If
\[
\mathcal T_n(q)=n!\prod_{m=1}^{n}(1-\alpha_m z),
\]
then
\[
p\boxplus_n q

=
 \Bigl(\prod_{m=1}^{n}(I-\alpha_m D)\Bigr)p,
\qquad D:=\frac{d}{dx}.
\]
Equivalently, if $q_{\alpha}(x):=x^n-n\alpha x^{n-1}$, then
\[
p\boxplus_n q
=
(\cdots((p\boxplus_n q_{\alpha_1})\boxplus_n q_{\alpha_2})\cdots)\boxplus_n q_{\alpha_n},
\]
and each elementary step satisfies
\[
p\boxplus_n q_{\alpha}=p-\alpha p'.
\]
\end{enumerate}
\end{lemma}

\begin{proof}
For (1), write
\[
p(x)=\sum_{i=0}^{n}a_i x^{n-i},\qquad
q(x)=\sum_{j=0}^{n}b_j x^{n-j},\qquad
r:=p\boxplus_n q=\sum_{k=0}^{n}c_k x^{n-k}.
\]
By definition,
\[
c_k=\sum_{i+j=k}\frac{(n-i)!(n-j)!}{n!(n-k)!}\,a_i b_j.
\]
Hence
\[
(n-k)!\,c_k=\frac1{n!}\sum_{i+j=k}(n-i)!\,a_i\,(n-j)!\,b_j,
\]
which is exactly coefficient-wise multiplication identity
\[
\mathcal T_n(r)=\frac1{n!}\,\mathcal T_n(p)\mathcal T_n(q).
\]

For (2), note first that for $q_{\alpha}(x)=x^n-n\alpha x^{n-1}$ we have
\[
\mathcal T_n(q_\alpha)=n!(1-\alpha z).
\]
Applying (1),
\[
\mathcal T_n(p\boxplus_n q_\alpha)=\frac1{n!}\,\mathcal T_n(p)\,n!(1-\alpha z)
=(1-\alpha z)\mathcal T_n(p).
\]
Since $\mathcal T_n$ is an isomorphism on degree-$n$ coefficient vectors, this realizes $p\boxplus_n q_\alpha$ as multiplication by $(1-\alpha z)$ in $\mathcal T_n$-space.

Now from Lemma~\ref{lem:q4-diffrepr},
\[
p\boxplus_n q_\alpha
=
\frac1{n!}\bigl(p\,q_\alpha^{(n)}(0)+p'\,q_\alpha^{(n-1)}(0)\bigr)
=
\frac1{n!}\bigl(p\,n!+p'(-\alpha n!)\bigr)=p-\alpha p'.
\]

Finally, because multiplication in $\mathbb R[z]/(z^{n+1})$ is associative and commutative,
\[
\frac1{n!}\mathcal T_n(p)\mathcal T_n(q)
=
\frac1{n!}\mathcal T_n(p)\,n!\prod_{m=1}^{n}(1-\alpha_m z)
=
\mathcal T_n(p)\prod_{m=1}^{n}(1-\alpha_m z),
\]
which equals the $\mathcal T_n$-image of successive convolutions by $q_{\alpha_m}$, i.e. successive application of $(I-\alpha_m D)$.
\end{proof}

\begin{lemma}[Derivative recursion for $\boxplus_n$]\label{lem:q4-deriv-recursion}
For monic degree-$n$ polynomials $p,q$,
\[
\frac{d}{dx}(p\boxplus_n q)
=
\frac1n\,(p'\boxplus_{n-1} q').
\]
More generally, for every integer $m$ with $0\le m\le n$,
\[
\frac{d^m}{dx^m}(p\boxplus_n q)
=
\frac{1}{n(n-1)\cdots(n-m+1)}
\bigl(p^{(m)}\boxplus_{n-m} q^{(m)}\bigr).
\]
\end{lemma}

\begin{proof}
Write
\[
p(x)=\sum_{i=0}^n a_i x^{n-i},\qquad
q(x)=\sum_{j=0}^n b_j x^{n-j},\qquad
r(x):=(p\boxplus_n q)(x)=\sum_{k=0}^n c_k x^{n-k}.
\]
By definition,
\[
c_k=\sum_{i+j=k}\frac{(n-i)!(n-j)!}{n!(n-k)!}\,a_i b_j.
\]
Hence
\[
r'(x)=\sum_{k=0}^{n-1}(n-k)c_k\,x^{n-1-k},
\]
so the coefficient of $x^{(n-1)-k}$ in $r'$ is
\[
(n-k)c_k
=
\sum_{i+j=k}\frac{(n-i)!(n-j)!}{n!(n-k-1)!}\,a_i b_j.
\tag{1}
\]

Now write
\[
p'(x)=\sum_{i=0}^{n-1}\alpha_i x^{(n-1)-i},
\quad
q'(x)=\sum_{j=0}^{n-1}\beta_j x^{(n-1)-j},
\]
with
\[
\alpha_i=(n-i)a_i,\qquad \beta_j=(n-j)b_j.
\]
In degree $n-1$, the $\boxplus_{n-1}$ coefficient formula gives:
the coefficient of $x^{(n-1)-k}$ in $(p'\boxplus_{n-1}q')$ equals
\[
\sum_{i+j=k}
\frac{(n-1-i)!(n-1-j)!}{(n-1)!(n-1-k)!}\,\alpha_i\beta_j
\]
\[
=
\sum_{i+j=k}
\frac{(n-i)!(n-j)!}{(n-1)!(n-k-1)!}\,a_i b_j.
\tag{2}
\]
Comparing (1) and (2), we get
\[
r'=\frac1n\,(p'\boxplus_{n-1}q').
\]
This proves the $m=1$ case.

Iterating the identity $m$ times yields
\[
\frac{d^m}{dx^m}(p\boxplus_n q)
=
\frac1n\cdot\frac1{n-1}\cdots\frac1{n-m+1}
\bigl(p^{(m)}\boxplus_{n-m} q^{(m)}\bigr),
\]
which is the stated general formula.
\end{proof}

\begin{lemma}[Heat-semigroup intertwining]\label{lem:q4-heat-intertwine}
Let $H_t:=e^{t\partial_{xx}}$ acting on degree-$n$ polynomials.
Then for all $t\in\mathbb R$,
\[
H_{2t}(p\boxplus_n q)=(H_t p)\boxplus_n(H_t q).
\]
\end{lemma}

\begin{proof}
Write
\[
p(x)=\sum_{i=0}^n a_i x^{n-i},\qquad
q(x)=\sum_{i=0}^n b_i x^{n-i}.
\]
Work in the truncated ring $\mathbb R[z]/(z^{n+1})$ and define
\[
\mathcal T(p):=\sum_{i=0}^n (n-i)!\,a_i z^i.
\]
From the coefficient formula of $\boxplus_n$,
\[
\mathcal T(p\boxplus_n q)=\frac1{n!}\,\mathcal T(p)\mathcal T(q)
\quad\text{in }\mathbb R[z]/(z^{n+1}).
\tag{1}
\]
Next, if $p''(x)=\sum_{i=0}^n a_i' x^{n-i}$, then
\[
a'_{i+2}=(n-i)(n-i-1)a_i
\]
for $0\le i\le n-2$, hence
\[
(n-(i+2))!\,a'_{i+2}=(n-i)!\,a_i.
\]
Therefore
\[
\mathcal T(p'')=z^2\mathcal T(p)
\quad\text{in }\mathbb R[z]/(z^{n+1}).
\tag{2}
\]
By linearity and power-series expansion, (2) gives
\[
\mathcal T(H_t p)=e^{t z^2}\mathcal T(p).
\tag{3}
\]
Now combine (1) and (3):
\[
\mathcal T\!\left((H_t p)\boxplus_n(H_t q)\right)
=\frac1{n!}\,e^{t z^2}\mathcal T(p)\,e^{t z^2}\mathcal T(q)
=\frac1{n!}\,e^{2t z^2}\mathcal T(p)\mathcal T(q)
=\mathcal T\!\left(H_{2t}(p\boxplus_n q)\right).
\]
Since $\mathcal T$ is injective on degree-$n$ coefficient vectors, the claimed identity follows.
\end{proof}

\begin{lemma}[Variational formula for $1/\Phi_n$]\label{lem:q4-variational}
If $p$ has simple real roots $\lambda_1,\dots,\lambda_n$, then
\[
\frac1{\Phi_n(p)}
=
\min_{\substack{w_{ij}\in\mathbb R\\ \sum_{i\neq j} w_{ij}=1}}
\sum_{i\neq j} w_{ij}^2(\lambda_i-\lambda_j)^2.
\]
The minimum is attained at
\[
w_{ij}=\frac{(\lambda_i-\lambda_j)^{-2}}{\sum_{a\neq b}(\lambda_a-\lambda_b)^{-2}}
=\frac{(\lambda_i-\lambda_j)^{-2}}{\Phi_n(p)}.
\]
\end{lemma}

\begin{proof}
Set $d_{ij}:=(\lambda_i-\lambda_j)^2>0$.
For any weights with $\sum_{i\neq j} w_{ij}=1$, Cauchy--Schwarz gives
\[
1^2
=\left(\sum_{i\neq j} w_{ij}\right)^2
=\left(\sum_{i\neq j} (w_{ij}\sqrt{d_{ij}})\,d_{ij}^{-1/2}\right)^2
\le
\left(\sum_{i\neq j} w_{ij}^2 d_{ij}\right)
\left(\sum_{i\neq j} d_{ij}^{-1}\right).
\]
Thus
\[
\sum_{i\neq j} w_{ij}^2 d_{ij}\ge \frac1{\sum_{i\neq j} d_{ij}^{-1}}
=\frac1{\Phi_n(p)}.
\]
Equality holds iff $w_{ij}\sqrt{d_{ij}}$ is proportional to $d_{ij}^{-1/2}$, i.e. $w_{ij}=c\,d_{ij}^{-1}$.
The constraint $\sum w_{ij}=1$ gives $c=1/\Phi_n(p)$.
\end{proof}

\begin{proposition}[Stam-type inequality for monotone root-sum model]\label{prop:q4-monotone-sum}
Let $\lambda_1\le\cdots\le\lambda_n$ and $\mu_1\le\cdots\le\mu_n$ be two real root vectors, and define
\[
\nu_i:=\lambda_i+\mu_i,\qquad i=1,\dots,n.
\]
Then
\[
\frac1{\Phi_n(\nu)}\ge \frac1{\Phi_n(\lambda)}+\frac1{\Phi_n(\mu)},
\]
where $\Phi_n(\xi):=\sum_{i\neq j}(\xi_i-\xi_j)^{-2}$ for simple configurations.
\end{proposition}

\begin{proof}
Let $\mathcal I:=\{(i,j):1\le i<j\le n\}$. For $\alpha=(i,j)\in\mathcal I$, define
\[
a_\alpha:=(\lambda_j-\lambda_i)^2,\qquad
b_\alpha:=(\mu_j-\mu_i)^2,\qquad
c_\alpha:=(\nu_j-\nu_i)^2.
\]
Since $\lambda,\mu$ are nondecreasing, all increments are nonnegative, so
\[
c_\alpha=( (\lambda_j-\lambda_i)+(\mu_j-\mu_i) )^2\ge a_\alpha+b_\alpha.
\]
Hence
\[
\sum_{\alpha\in\mathcal I}\frac1{c_\alpha}\le \sum_{\alpha\in\mathcal I}\frac1{a_\alpha+b_\alpha}.
\]
Define for any positive vector $t=(t_\alpha)_{\alpha\in\mathcal I}$:
\[
H(t):=\frac1{\sum_{\alpha\in\mathcal I} t_\alpha^{-1}}.
\]
Then the previous inequality is $H(c)\ge H(a+b)$.

We now show $H$ is superadditive:
\[
H(u+v)\ge H(u)+H(v)\quad\text{for all positive }u,v.
\]
By the same Cauchy--Schwarz argument as Lemma~\ref{lem:q4-variational},
\[
H(t)=\min_{\substack{\theta_\alpha\in\mathbb R\\ \sum_\alpha \theta_\alpha=1}}
\sum_\alpha \theta_\alpha^2 t_\alpha.
\]
Therefore
\[
H(u+v)
=\min_\theta \sum_\alpha \theta_\alpha^2(u_\alpha+v_\alpha)
\ge
\min_\theta \sum_\alpha \theta_\alpha^2u_\alpha
+
\min_\theta \sum_\alpha \theta_\alpha^2v_\alpha
=H(u)+H(v).
\]
Applying this with $u=a$, $v=b$ gives
\[
H(c)\ge H(a+b)\ge H(a)+H(b).
\]
Finally,
\[
\frac1{\Phi_n(\lambda)}=\frac{H(a)}{2},\quad
\frac1{\Phi_n(\mu)}=\frac{H(b)}{2},\quad
\frac1{\Phi_n(\nu)}=\frac{H(c)}{2},
\]
because $\Phi_n(\xi)=2\sum_{i<j}(\xi_j-\xi_i)^{-2}$.
Thus
\[
\frac1{\Phi_n(\nu)}\ge \frac1{\Phi_n(\lambda)}+\frac1{\Phi_n(\mu)}.
\]
\end{proof}

\begin{remark}[Scope of Proposition~\ref{prop:q4-monotone-sum}]
Proposition~\ref{prop:q4-monotone-sum} is a full all-$n$ result for the deterministic monotone coupling model
$\nu_i=\lambda_i+\mu_i$.
It does not by itself prove \eqref{eq:q4-star}, because in general the root vector of $p\boxplus_n q$ is not equal to this monotone sum vector.
\end{remark}

%% ---- Semi-Gaussian Stam inequality (all n) ----

\begin{lemma}[Edge-score first-moment identity]\label{lem:q4-edge-score}
For $p$ with simple roots $\lambda_1<\cdots<\lambda_n$ and edge scores
$g_{ij}:=(s_i-s_j)/(\lambda_i-\lambda_j)$, $1\le i<j\le n$:
\[
\sum_{i<j}g_{ij}=\Phi_n(p).
\]
\end{lemma}

\begin{proof}
$\Phi_n=\sum_i s_i^2
=\sum_i s_i\sum_{j\ne i}(\lambda_i-\lambda_j)^{-1}
=\sum_{i<j}\bigl(s_i(\lambda_i-\lambda_j)^{-1}+s_j(\lambda_j-\lambda_i)^{-1}\bigr)
=\sum_{i<j}\frac{s_i-s_j}{\lambda_i-\lambda_j}
=\sum_{i<j}g_{ij}$.
\end{proof}

\begin{lemma}[Fisher information of Hermite polynomials]\label{lem:q4-hermite-phi}
For the probabilist's Hermite polynomial
$\mathrm{He}_n(x)=x^n-\tbinom{n}{2}x^{n-2}+\cdots$
with simple real roots $\lambda_1<\cdots<\lambda_n$:
\[
\Phi_n(\mathrm{He}_n)=\frac{n(n-1)}{4}.
\]
For a scaled Hermite polynomial with roots $\sqrt{s}\,\lambda_i$
(i.e., the polynomial $s^{n/2}\mathrm{He}_n(x/\sqrt{s})$):
\[
\Phi_n(\sqrt{s}\,\mathrm{He}_n)=\frac{n(n-1)}{4s},
\qquad
\frac{1}{\Phi_n(\sqrt{s}\,\mathrm{He}_n)}=\frac{4s}{n(n-1)}.
\]
\end{lemma}

\begin{proof}
The Hermite differential equation
$\mathrm{He}_n''(x)-x\,\mathrm{He}_n'(x)+n\,\mathrm{He}_n(x)=0$
evaluated at a root $\lambda_i$ (where $\mathrm{He}_n(\lambda_i)=0$) gives
$\mathrm{He}_n''(\lambda_i)=\lambda_i\,\mathrm{He}_n'(\lambda_i)$.
By the standard identity $p''(\lambda_i)/p'(\lambda_i)=2s_i$,
we get $s_i=\lambda_i/2$.
Hence
\[
\Phi_n(\mathrm{He}_n)
=\sum_{i=1}^n s_i^2
=\frac{1}{4}\sum_{i=1}^n\lambda_i^2.
\]
By Vieta's formulas for $\mathrm{He}_n$:
$\sum_i\lambda_i=0$ and $\sum_{i<j}\lambda_i\lambda_j=-\binom{n}{2}$,
so
$\sum_i\lambda_i^2=(\sum_i\lambda_i)^2-2\sum_{i<j}\lambda_i\lambda_j=n(n-1)$.

The scaling claim follows from $\Phi_n(c\cdot p(\cdot/c))=\Phi_n(p)/c^2$,
which is immediate from the definition (root gaps scale by $c$).
\end{proof}

\begin{lemma}[Scaled-Hermite transform in $\mathcal T_n$]\label{lem:q4-hermite-transform}
Let $\mathrm{He}_n$ be the probabilist's Hermite polynomial, and define
\[
q_t(x):=t^{n/2}\,\mathrm{He}_n(x/\sqrt t)\qquad (t>0).
\]
Then in $\mathbb R[z]/(z^{n+1})$,
\[
\mathcal T_n(q_t)=n!\,e^{-t z^2/2}.
\]
\end{lemma}

\begin{proof}
The explicit coefficient formula is
\[
\mathrm{He}_n(x)=n!\sum_{m=0}^{\lfloor n/2\rfloor}
\frac{(-1)^m}{2^m m!(n-2m)!}\,x^{n-2m}.
\]
Hence
\[
q_t(x)
=n!\sum_{m=0}^{\lfloor n/2\rfloor}
\frac{(-1)^m t^m}{2^m m!(n-2m)!}\,x^{n-2m}.
\]
So the coefficient of $x^{n-2m}$ equals
$n!(-1)^m t^m/(2^m m!(n-2m)!)$.
Applying $\mathcal T_n$ multiplies this by $(n-2m)!$, yielding
\[
\mathcal T_n(q_t)
=n!\sum_{m=0}^{\lfloor n/2\rfloor}\frac{(-t z^2/2)^m}{m!}
=n!\,e^{-t z^2/2}
\quad\text{in }\mathbb R[z]/(z^{n+1}).
\]
\end{proof}

\begin{lemma}[Semi-Gaussian convolution is backward heat]\label{lem:q4-semi-gauss-flow}
For every degree-$n$ polynomial $p$ and every $t\ge 0$,
\[
p\boxplus_n q_t = H_{-t/2}p,
\qquad q_t(x):=t^{n/2}\mathrm{He}_n(x/\sqrt t).
\]
\end{lemma}

\begin{proof}
By Lemma~\ref{lem:q4-transform-factor},
\[
\mathcal T_n(p\boxplus_n q_t)=\frac1{n!}\mathcal T_n(p)\mathcal T_n(q_t).
\]
By Lemma~\ref{lem:q4-hermite-transform},
\[
\mathcal T_n(p\boxplus_n q_t)=\mathcal T_n(p)\,e^{-t z^2/2}.
\]
By Lemma~\ref{lem:q4-heat-intertwine} (formula $\mathcal T(H_s p)=e^{s z^2}\mathcal T(p)$ in its proof),
\[
\mathcal T_n(H_{-t/2}p)=e^{-t z^2/2}\mathcal T_n(p).
\]
Since $\mathcal T_n$ is injective on degree-$n$ coefficient vectors,
$p\boxplus_n q_t=H_{-t/2}p$.
\end{proof}

\begin{theorem}[Semi-Gaussian Stam inequality---all $n$]\label{thm:q4-semi-gaussian}
For every $n\ge 2$, every monic real-rooted polynomial $p$ of degree~$n$,
and every $s>0$, assume that the path
\[
p_t:=p\boxplus_n \sqrt t\,\mathrm{He}_n,\qquad t\in[0,s],
\]
has simple real roots for all $t\in(0,s]$. Then
\[
\frac{1}{\Phi_n(p\boxplus_n\sqrt{s}\,\mathrm{He}_n)}
\;\ge\;
\frac{1}{\Phi_n(p)}+\frac{1}{\Phi_n(\sqrt{s}\,\mathrm{He}_n)}.
\]
Equality holds for the Hermite family
$p(x)=a^{n/2}\mathrm{He}_n(x/\sqrt a)$ ($a>0$).
\end{theorem}

\begin{proof}
Set $q_t(x):=t^{n/2}\mathrm{He}_n(x/\sqrt t)$ and
\[
p_t:=p\boxplus_n q_t.
\]
By Lemma~\ref{lem:q4-semi-gauss-flow}, $p_t=H_{-t/2}p$, hence
\[
\partial_t p_t=-\frac12\,\partial_{xx}p_t.
\]
Let $\lambda_1(t)<\cdots<\lambda_n(t)$ be the roots of $p_t$.
Differentiating $p_t(\lambda_i(t))=0$ gives
\[
0=\partial_t p_t(\lambda_i)+p_t'(\lambda_i)\dot\lambda_i
=-\frac12 p_t''(\lambda_i)+p_t'(\lambda_i)\dot\lambda_i.
\]
Using $p_t''(\lambda_i)/p_t'(\lambda_i)=2s_i$ with
$s_i:=\sum_{j\neq i}(\lambda_i-\lambda_j)^{-1}$, we obtain
\[
\dot\lambda_i=s_i.
\]

\textbf{Step 1 (evolution of $\Phi_n$).}
Differentiate $s_i=\sum_{j\neq i}(\lambda_i-\lambda_j)^{-1}$:
\[
\dot s_i=\sum_{j\neq i}\frac{\dot\lambda_j-\dot\lambda_i}{(\lambda_i-\lambda_j)^2}.
\]
Therefore
\[
\dot\Phi_n
=2\sum_i s_i\dot s_i
=2\sum_{i<j}\frac{(s_i-s_j)(\dot\lambda_j-\dot\lambda_i)}{(\lambda_i-\lambda_j)^2}
=-2\sum_{i<j}\frac{(s_i-s_j)^2}{(\lambda_i-\lambda_j)^2}.
\]
With $g_{ij}:=(s_i-s_j)/(\lambda_i-\lambda_j)$, this is
\[
\dot\Phi_n=-2\sum_{i<j}g_{ij}^2\le 0,
\]
so for $J(t):=1/\Phi_n(p_t)$,
\[
J'(t)=\frac{-\dot\Phi_n}{\Phi_n^2}
=\frac{2\sum_{i<j}g_{ij}^2}{\Phi_n^2}.
\]

\textbf{Step 2 (Cauchy--Schwarz lower bound on $J'$).}
By Lemma~\ref{lem:q4-edge-score}: $\Phi_n=\sum_{i<j}g_{ij}$.
Cauchy--Schwarz over the $\binom{n}{2}$ pairs gives
\[
\Bigl(\sum_{i<j}g_{ij}\Bigr)^{\!2}
\;\le\; \binom{n}{2}\sum_{i<j}g_{ij}^2,
\]
hence
\[
J'(t)
\;\ge\; \frac{2}{\binom{n}{2}}
= \frac{4}{n(n-1)}
= \frac{1}{\Phi_n(\mathrm{He}_n)}
\qquad\text{(by Lemma~\ref{lem:q4-hermite-phi})}.
\]

\textbf{Step 3 (integration).}
Integrating from $0$ to $s$:
\[
J(s)-J(0)
\;\ge\; \frac{s}{\Phi_n(\mathrm{He}_n)}
= \frac{4s}{n(n-1)}
= \frac{1}{\Phi_n(\sqrt{s}\,\mathrm{He}_n)},
\]
which is exactly $1/\Phi_n(p_s)\ge 1/\Phi_n(p)+1/\Phi_n(\sqrt{s}\,\mathrm{He}_n)$.

\textbf{Hermite equality family.}
If $p(x)=a^{n/2}\mathrm{He}_n(x/\sqrt a)$, then
$\mathcal T_n(p)=n!e^{-a z^2/2}$.
Hence
\[
\mathcal T_n(p_t)=\frac1{n!}\mathcal T_n(p)\mathcal T_n(q_t)
=n!e^{-(a+t)z^2/2},
\]
so $p_t(x)=(a+t)^{n/2}\mathrm{He}_n(x/\sqrt{a+t})$.
By Lemma~\ref{lem:q4-hermite-phi},
\[
\frac{1}{\Phi_n(p_t)}=\frac{4(a+t)}{n(n-1)},
\]
which is affine in $t$ with slope $4/(n(n-1))$; therefore every inequality
step above is an equality for this family.
\end{proof}

\begin{remark}[Significance of the semi-Gaussian inequality]
Theorem~\ref{thm:q4-semi-gaussian} establishes
the target inequality~\eqref{eq:q4-star} whenever one of the two
polynomials is a scaled Hermite polynomial---for \textbf{all degrees} $n$.
This is the polynomial analogue of the classical result that the
Stam inequality holds when one summand is Gaussian, and it strictly
extends the degenerate shift case (Proposition~\ref{prop:q4-shift},
which requires one polynomial to be $(x-a)^n$).
Together with Propositions~\ref{prop:q4-n2} and~\ref{prop:q4-n3}
(which handle all $p,q$ for $n\le 3$) and
Proposition~\ref{prop:q4-monotone-sum}
(which handles the monotone coupling for all $n$),
this establishes~\eqref{eq:q4-star} in three independent
infinite families of cases.
\end{remark}

\begin{remark}[Status of equality classification]
The proof above gives a complete inequality proof under the stated
regularity assumption and a full explicit equality family (scaled Hermite).
A global characterization of all equality cases for this semi-Gaussian flow
is not used elsewhere in this paper.
\end{remark}

%% ---- End of Semi-Gaussian block ----

\begin{proposition}[Reduction of \eqref{eq:q4-star} to a weighted quadratic inequality]\label{prop:q4-reductionA}
For roots $\lambda$ of $p$, $\mu$ of $q$, and $\gamma$ of $p\boxplus_n q$, define
\[
\mathcal Q_w(\xi):=\sum_{i\neq j} w_{ij}^2(\xi_i-\xi_j)^2
\quad\text{for}\quad
w=(w_{ij})_{i\neq j},\ \sum_{i\neq j}w_{ij}=1.
\]
Assume the following statement holds for all such $w$:
\[
\mathcal Q_w(\gamma)\ge \mathcal Q_w(\lambda)+\mathcal Q_w(\mu).
\tag{A}\label{eq:q4-A}
\]
Then \eqref{eq:q4-star} follows.
\end{proposition}

\begin{proof}
By Lemma~\ref{lem:q4-variational},
\[
\frac1{\Phi_n(p\boxplus_n q)}=\min_w \mathcal Q_w(\gamma).
\]
If \eqref{eq:q4-A} holds for every $w$, then
\[
\min_w \mathcal Q_w(\gamma)
\ge
\min_w\big(\mathcal Q_w(\lambda)+\mathcal Q_w(\mu)\big)
\ge
\min_w \mathcal Q_w(\lambda)+\min_w \mathcal Q_w(\mu),
\]
and another application of Lemma~\ref{lem:q4-variational} gives
\[
\frac1{\Phi_n(p\boxplus_n q)}
\ge
\frac1{\Phi_n(p)}+\frac1{\Phi_n(q)}.
\]
\end{proof}

\begin{remark}[General-$n$ bottleneck in one line]
The all-$n$ closure of \eqref{eq:q4-star} is reduced to proving \eqref{eq:q4-A} uniformly in $n$ and $w$.
This isolates the remaining difficulty to a single weighted quadratic superadditivity statement at the root level.
\end{remark}

\begin{remark}[Important correction: \eqref{eq:q4-A} is sufficient but not necessary]
Proposition~\ref{prop:q4-reductionA} is logically correct as a \emph{sufficient} criterion, but the pointwise inequality \eqref{eq:q4-A} is generally too strong.
Numerical experiments for $n=4$ produce triples $(p,q,w)$ with
\[
\mathcal Q_w(\gamma)<\mathcal Q_w(\lambda)+\mathcal Q_w(\mu),
\]
while the target Stam-type inequality \eqref{eq:q4-star} still holds on the same $(p,q)$.
Therefore the proof of \eqref{eq:q4-star} must use a weaker global mechanism than uniform-in-$w$ superadditivity.
\end{remark}

\begin{remark}[Program toward the full theorem]
A viable closure route for \eqref{eq:q4-star} is to prove a finite-degree entropy-power inequality for a discriminant-based functional compatible with $\boxplus_n$. Combined with Lemma~\ref{lem:q4-debruijn}, this would mirror the classical entropy-power $\Rightarrow$ Stam mechanism.
\end{remark}

\begin{remark}[Additional stress-tests after rejecting \eqref{eq:q4-A}]
After identifying \eqref{eq:q4-A} as too strong, we ran three independent numerical checks:
\begin{enumerate}
\item Random search: for $n=4,5,6,7,8$, $6\times 10^4$ random pairs $(p,q)$ per $n$ (all real-rooted inputs) produced no violation of \eqref{eq:q4-star}; observed margins stayed positive.
\item Local descent: adversarial local optimization for $n=4$ drives the margin close to $0$ but high-precision recomputation stays nonnegative.
\item Near-degenerate regime: with one input approaching $(x-a)^n$, the margin scales as $+C\varepsilon^2+o(\varepsilon^2)$ in tested families, consistent with stability of the equality case in Proposition~\ref{prop:q4-shift}.
\end{enumerate}
These tests strengthen confidence in \eqref{eq:q4-star}, but do not constitute a proof.
\end{remark}

\begin{remark}[Current conclusion]
Within this draft, \eqref{eq:q4-star} is proved in the following cases:
\begin{enumerate}[nosep]
\item $n=2$, all $p,q$: exact equality (Proposition~\ref{prop:q4-n2}).
\item $n=3$, all $p,q$: strict inequality (Proposition~\ref{prop:q4-n3}).
\item All $n$, $p=(x-a)^n$: equality (Proposition~\ref{prop:q4-shift}).
\item All $n$, monotone coupling $\nu_i=\lambda_i+\mu_i$
  (Proposition~\ref{prop:q4-monotone-sum}).
\item \textbf{All $n$, $q=\sqrt{s}\,\mathrm{He}_n$}
  (Theorem~\ref{thm:q4-semi-gaussian}, under its regularity hypothesis):
  the semi-Gaussian case.
\end{enumerate}
The general all-$n$ statement for arbitrary $p,q$ with $n\ge 4$ remains open.
\end{remark}
